\documentclass[11pt]{article}

\usepackage{aa222-jmlr2e}
\makeatother
\usepackage{lipsum}
\usepackage{pdfpages}
\usepackage{sectsty}
\usepackage{hyperref}
\usepackage{graphicx}
\usepackage{subcaption}
\usepackage{float}
\usepackage{amsmath}
\usepackage{booktabs}
\usepackage{natbib}
\bibliographystyle{ieeetr}

\begin{document}

\title{CME 213 Final Project Milestone 3\\
Real-Time Multi-Camera 3D Perception and Tracking for Robotics}

\author{%
\begin{center}
\begin{minipage}[t]{0.46\textwidth}\centering
\textbf{Frances Raphael}\\
{\ttfamily fraphael@stanford.edu}\\
\end{minipage}\hfill
\begin{minipage}[t]{0.46\textwidth}\centering
\textbf{Aanya Tashfeen}\\
{\ttfamily aanya@stanford.edu}\\
\end{minipage}
\end{center}
}
\maketitle

\section{Summary of Progress}

Since Milestone 2, we have completed the CPU baseline and all three planned CUDA kernel variants for
the stereo SAD disparity estimation step of our pipeline. The CPU baseline described in our Milestone 2
design is fully implemented and verified. The three GPU kernels---a naive global-memory kernel, a
shared-memory tiled kernel, and a row-reuse tiled kernel---are all implemented, launched from a
common host driver, and produce verified correct results. A benchmarking driver and CSV logging
infrastructure are also in place to facilitate the performance analysis planned for Milestone 4.

The main change from our Milestone 2 design is that we narrowed the initial GPU target to a single
kernel: stereo SAD matching. In Milestone 2 we left open the possibility of also accelerating
preprocessing, but we found that the SAD matching step dominates runtime so completely that
it is the only kernel worth optimizing at this stage. Object tracking and MPI distribution are not
yet started, as planned for Milestone 4.

\section{CUDA Kernels}

We implemented three CUDA kernels, all computing the same operation: for each pixel $(r, c)$ in
the left image, find the disparity $d \in [0, D_{\max})$ that minimizes the Sum of Absolute Differences
(SAD) over a $(2R+1) \times (2R+1)$ patch between the left image centred at $(r,c)$ and the right
image centred at $(r, c-d)$. All kernels use a one-thread-per-output-pixel decomposition with
$16 \times 16$ or $32 \times 8$ thread blocks.

\textbf{Basic kernel.} Each thread independently fetches all required pixels from global memory.
For every candidate disparity, the thread reads a $(2R+1)^2$ patch from each image and accumulates
the absolute differences. No data is shared between threads. This kernel serves as the lower-bound
baseline for the GPU implementations and stresses global memory bandwidth, since each pixel
may be loaded by many overlapping patches across nearby threads.

\textbf{Shared memory kernel.} Each $16 \times 16$ thread block cooperatively loads a tile of the left
image and a wider tile of the right image into shared memory before any SAD computation begins.
The right image tile is extended by $D_{\max}$ columns to cover every disparity offset needed by
threads in the block. After a \texttt{\_\_syncthreads()} barrier, each thread reads its patch
entirely from on-chip shared memory. This eliminates the redundant global memory reads that occur
in the basic kernel when adjacent threads access overlapping patches, and is expected to be the
most memory-efficient variant.

\textbf{Tiled kernel.} This variant eliminates register spilling by declaring the accumulator
array with a compile-time size \texttt{TILED\_MAX\_DISP}$= 64$ and applying \texttt{\#pragma unroll}
to all disparity loops. After unrolling, each \texttt{sad[d]} has a statically known index,
so the compiler assigns it its own register rather than spilling to local memory.
The patch is traversed once; for each position the left-image pixel is loaded into a register
and reused across all 64 disparity comparisons, giving a $64\times$ reduction in left-image
loads relative to the basic kernel. Right-image accesses are a descending consecutive sequence
per patch element per warp, which coalesces into a single cache-line transaction.
The kernel uses $16 \times 8$ thread blocks, giving 128 threads per block and keeping
total register usage (64 accumulators $+$ overhead $\approx$ 80 registers per thread $\times$
128 threads $= 10{,}240$ registers) well within the 65{,}536-register SM budget.

\section{C++ Wrapper and CPU Code}

All GPU memory management is handled by a \texttt{GpuBuffers} RAII struct that allocates
device memory for the left image, right image, and disparity output on construction, copies the
host images to the device via \texttt{cudaMemcpy}, and frees all allocations on destruction.
Kernel launches are wrapped in a common timing helper that uses \texttt{cudaEvent\_t} to measure
elapsed time with sub-millisecond precision. Each kernel is warmed up with one untimed launch
before the timed repetitions begin, to avoid including JIT compilation or cold-start overhead
in the reported measurements.

The CPU baseline implements the same SAD algorithm in plain C++ using nested loops. It serves
as both a performance reference and a correctness oracle for all three GPU kernels. A
\texttt{time\_sad\_stereo\_cpu} function runs the baseline for a configurable number of repetitions
and reports mean, standard deviation, and minimum wall-clock time. A shared benchmarking driver
(\texttt{main\_gpu}) runs all four implementations in sequence on the same synthetic image pair
and logs one CSV row per implementation, which is used to generate roofline plots offline.
Image outputs (left image, right image, disparity maps) are saved as binary PGM files for
visual inspection.

\section{Correctness Testing}

We verify each GPU kernel against the CPU reference. For a synthetic stereo pair where the right
image is the left image shifted by a known disparity $d^* = 24$ pixels, the minimum SAD at
$d = d^*$ is identically zero, which is unbeatable. Any correct implementation should therefore
recover $d^*$ for every valid pixel, giving mean absolute error (MAE) of exactly zero.

An early run revealed a mismatch between the ground truth validity mask (which marked pixels
as valid when $c \geq d^* + R$) and the SAD kernel's validity check (which requires $c \geq D_{\max}
- 1 + R$ to ensure all disparities are in bounds). This caused $\sim$6.4\% of evaluated pixels
to receive the border fill value of 0 rather than the correct disparity, inflating MAE to 1.53 px.
After aligning the ground truth mask with the kernel's validity condition, all four implementations---CPU
and all three GPU kernels---achieve MAE~$= 0.000$~px and bad-pixel rate~$= 0.00\%$ on the
480~$\times$~640 test image with $D_{\max} = 64$, $R = 2$.

\section{Preliminary Performance Analysis}

Table~\ref{tab:perf} reports measured runtimes on an NVIDIA RTX 2080 Ti (Turing, sm\_75,
616~GB/s memory bandwidth, $\sim$13.4~TOPS INT32 throughput) for a 480~$\times$~640 image
with $D_{\max} = 64$, $R = 2$ (5~$\times$~5 patches, approximately 0.87~GOps per frame).
Each measurement is the mean over five timed repetitions after one warm-up run.

\begin{table}[H]
\centering
\caption{Performance comparison across implementations. Speedup is relative to the CPU baseline.
All GPU results have MAE $= 0.000$~px against the CPU reference.}
\label{tab:perf}
\begin{tabular}{lrrrr}
\toprule
Implementation & Time (ms) & FPS & Throughput (GOps/s) & Speedup \\
\midrule
CPU (single-thread) & 359.4 & 2.8     & 2.4     & 1$\times$   \\
GPU basic           & 0.842 & 1{,}187 & 1{,}034 & 427$\times$ \\
GPU shared memory   & 0.414 & 2{,}417 & 2{,}106 & 868$\times$ \\
GPU tiled           & 0.429 & 2{,}330 & 2{,}030 & 838$\times$ \\
\bottomrule
\end{tabular}
\end{table}

The shared memory and tiled kernels achieve nearly identical performance---0.414~ms and 0.429~ms
respectively---both delivering roughly $868\times$ and $838\times$ speedup over the CPU and
approximately 2{,}100~GOps/s throughput, corresponding to $\sim$15\% of the GPU's theoretical
INT32 peak. The basic kernel is $\sim$2.5$\times$ slower at 0.842~ms, limited by redundant global
memory traffic from overlapping patches across adjacent threads.

Because the 480~$\times$~640 input images ($\approx$0.6~MB total) fit entirely in the RTX
2080 Ti's 6~MB L2 cache, even the basic kernel exceeds the 616~GB/s DRAM bandwidth ceiling:
its apparent throughput of 1{,}034~GOps/s implies an effective bandwidth of $\sim$1{,}034~GB/s,
only possible because repeated accesses hit L2 rather than DRAM. The smem and tiled kernels
push this further by reducing total data movement. The smem kernel loads each image tile once
per block and eliminates inter-thread redundancy in shared memory. The tiled kernel achieves
the same result differently: by loading each left-image pixel once into a register and reusing
it across all 64 disparity comparisons, it cuts left-image loads by $64\times$ relative to the
basic kernel. Both strategies converge on the same throughput, suggesting the remaining
bottleneck is compute throughput rather than memory bandwidth at this image size.

An earlier version of the tiled kernel suffered register spilling: allocating \texttt{sad[max\_disp]}
with a runtime bound caused the compiler to place the array in local memory (backed by DRAM),
degrading runtime to 7.78~ms---worse than basic. The fix was to declare the array with the
compile-time constant \texttt{TILED\_MAX\_DISP}$=$64 and apply \texttt{\#pragma unroll} to all
disparity loops, giving each \texttt{sad[d]} its own register after unrolling and eliminating
all local memory traffic.

\section{Challenges and Next Steps}

Two main technical challenges arose during this milestone. The first was the ground truth
validity mismatch described in Section~4. The second was diagnosing and fixing the tiled
kernel's register spilling: our throughput metric counts arithmetic operations, so a kernel
that spills to local memory simply appears to do less useful work per second rather than
explicitly revealing excess memory traffic. Identifying the root cause required reasoning
about compiler register allocation and reconfirming with measured runtime (7.78~ms with
spilling vs.\ 0.429~ms after the fix). Both issues are now resolved, with all four
implementations producing MAE $= 0.000$~px.

For Milestone 4, the remaining tasks are:

\textbf{MPI distribution.} Partition the image rows across MPI processes, with each process
running the shared-memory GPU kernel on its assigned slab. Implement halo exchange to handle
the patch border at slab boundaries, where each process needs \texttt{radius} rows from its
neighbors to compute a complete disparity map.

\textbf{Larger image experiments.} Benchmark on 1080p images ($1920 \times 1080$) and with
larger $D_{\max}$ to test whether the L2 caching benefit observed at 480~$\times$~640
degrades when the image data exceeds the 6~MB L2 cache.

\textbf{Strong scaling.} Measure runtime and efficiency as the number of MPI processes increases
from 1 to 8 with the total image size held fixed, and report communication overhead relative
to GPU kernel time.

\newpage

\bibliography{refs}

\end{document}
